\begin{mistake}{2026-04-09}{17}

\begin{problemcontent}
设 \( f(x) \) 与 \( g(x) \) 在 \( (-\infty, +\infty) \) 内都有定义，且 \( x=x_1 \) 是 \( f(x) \) 的唯一间断点，\( x=x_2 \) 是 \( g(x) \) 的唯一间断点，则：

(A) 当 \( x_1=x_2 \) 时，\( f(x)+g(x) \) 必有唯一的间断点 \( x=x_1 \)。

(B) 当 \( x_1 \neq x_2 \) 时，\( f(x)+g(x) \) 必有两个间断点 \( x=x_1 \) 与 \( x=x_2 \)。

(C) 当 \( x_1=x_2 \) 时，\( f(x)g(x) \) 必有唯一的间断点 \( x=x_1 \)。

(D) 当 \( x_1 \neq x_2 \) 时，\( f(x)g(x) \) 必有两个间断点 \( x=x_1 \) 与 \( x=x_2 \)。
\end{problemcontent}

\begin{solutioncontent}
正确答案是 \textbf{B}。

分析 \(F(x) = f(x) + g(x)\) 的连续性：
已知 \(x_1 \neq x_2\)，考察点 \(x_1\)：
在 \(x = x_1\) 处，\(f(x)\) 是间断的，而 \(x_1\) 不是 \(g(x)\) 的间断点（因为 \(g(x)\) 的唯一间断点是 \(x_2\)），所以 \(g(x)\) 在 \(x_1\) 处连续。
由连续函数的性质：\textbf{间断函数 \(\pm\) 连续函数 = 必间断}。
所以 \(F(x)\) 在 \(x_1\) 处必间断。
同理，在 \(x_2\) 处，\(f(x)\) 连续而 \(g(x)\) 间断，\(F(x)\) 在 \(x_2\) 处必间断。
其他所有点处两函数均连续，和必连续。故有且仅有两个间断点。

对于A、C、D项的反例：
若 \(x_1 = x_2\)，可能存在互补使和连续（如 \(f(x)\) 在 \(x_1\) 处跳跃，\(g(x)=-f(x)\)，和为常数连续），故A错。
对于乘积，\textbf{间断函数 \(\times\) 连续函数} 的结果不一定间断（如果连续函数在该点函数值为0，可能把间断给“抹平”），故C、D错。
\end{solutioncontent}

\mistakeknowledge{
\begin{itemize}
 \item \textbf{核心公式/定理}：连续性四则运算法则：连续\(\pm\)间断=必间断；间断\(\pm\)间断=不一定。
 \item \textbf{易错点}：误以为两个间断点相异时，乘积 \(f(x)g(x)\) 也一定在两处间断。若连续因子在该点为 \(0\)，则可能将无界或跳跃抵消（例如 \(f(x)=1/x\), \(g(x)=x^2\)，在 \(x=0\) 处 \(f(x)g(x)=x\) 变连续，补充定义后）。
 \item \textbf{关联知识}：“连续\(\times\)间断”的保号性推论。若要判定“连续\(\times\)间断=间断”，需满足连续函数在该点的值 \(\neq 0\)。
\end{itemize}}

\end{mistake}
